\chapter{Đọc thêm 2: Tác động của việc đánh cắp dữ liệu lên giá trị thị trường của doanh nghiệp}
\label{ch:M2L4-reading2}

% Nguồn: Chang, Kuo-Chung. "The Effect of Data Theft on a Firm's Short-Term and Long-Term Market Value." \emph{Mathematics}, vol. 8, no. 5, 2020. https://www.mdpi.com/2227-7390/8/5/808/htm
% Phần cần đọc: tr. 6--7 (~10 phút)
% Nội dung bắt buộc đọc: được kiểm tra trong mọi bài đánh giá (kể cả Quiz).

\section{Thông tin nguồn}

\begin{itemize}
  \item \textbf{Nguồn:} Chang, Kuo-Chung. "The Effect of Data Theft on a Firm's Short-Term and Long-Term Market Value." \emph{Mathematics}, vol. 8, no. 5, 2020. \url{https://www.mdpi.com/2227-7390/8/5/808/htm}
  \item \textbf{Phần cần đọc:} tr. 6--7 (~10 phút)
\end{itemize}

% [PDF p.6 | in p.6]
\section{Phát triển các Giả thuyết}

Các vụ vi phạm dữ liệu (data breaches) khiến tổ chức phải chịu thêm chi phí nhân công để khắc phục sự cố và sửa chữa các dữ liệu và hệ thống bị hỏng. Chúng cũng có thể dẫn đến việc giảm năng suất và giảm doanh thu do thời gian ngừng hoạt động (downtime) ngoài dự kiến. Các chi phí bổ sung khác bao gồm thiết lập đường dây nóng cho khách hàng, cung cấp dịch vụ giám sát tín dụng cho nạn nhân cũng như các khoản phí kiện tụng. Những chi phí này có thể tác động tạm thời đến hoạt động kinh doanh hoặc dẫn đến những khoản phí tức thì đối với khả năng sinh lời dự kiến của công ty. Do đó, một công ty bị vi phạm dữ liệu dự kiến sẽ gặp phải sự sụt giảm trong dòng tiền thuần tương lai của mình. Kết quả là, các nhà đầu tư sẽ đánh giá lại định giá của họ đối với công ty. Vì vụ vi phạm dự kiến sẽ có tác động tiêu cực đến dòng tiền thuần, các mức định giá cũng dự kiến sẽ giảm xuống. Như vậy, chúng tôi đưa ra giả thuyết:

\textbf{Giả thuyết 1 (H1)}. \emph{Việc công bố một vụ vi phạm dữ liệu có tác động tiêu cực đến giá trị thị trường ngắn hạn của công ty bị vi phạm.}

Ngoài những hậu quả mang tính nhất thời đối với hiệu quả tài chính, tác động của các vụ vi phạm dữ liệu là rất sâu rộng và có thể gây tổn hại cho các công ty ở bất kỳ quy mô nào. Những hệ lụy của các sự cố vi phạm có thể tiếp diễn trong nhiều năm tới. Bên cạnh các chi phí trả trước cho việc thông báo tới các chủ thể dữ liệu cùng với việc điều tra và kiểm soát vụ vi phạm, các tập đoàn cũng phải đối mặt với nguy cơ bị kiện tụng và nộp phạt, cũng như các chi phí vô hình khác liên quan đến thiệt hại về thương hiệu và uy tín doanh nghiệp. Sự sút giảm uy tín này có thể đồng nghĩa với việc mất khách hàng cũng như sự sụt giảm trong việc thu hút khách hàng mới, bởi vì mối quan hệ giữa an ninh dữ liệu và định giá thị trường của công ty có thể được quy cho niềm tin của khách hàng. Các nghiên cứu đã chỉ ra rằng an ninh dữ liệu là thiết yếu đối với niềm tin của khách hàng do những lo ngại liên quan đến quyền riêng tư thông tin [38--40]. Những khách hàng hiện tại có thể không sẵn lòng giao dịch với công ty bị vi phạm vì những lo ngại về quyền riêng tư. Hậu quả là, công ty bị vi phạm phải chịu sự sụt giảm lòng trung thành của khách hàng cũng như tốc độ thu hút khách hàng mới bị thu hẹp. Hơn nữa, các vụ vi phạm an ninh cũng có thể dẫn đến chi phí kinh doanh trong tương lai cao hơn do sự e ngại của các đối tác hoặc nhà cung cấp khi tiếp tục quan hệ hợp tác với các công ty bị vi phạm. Từ đó, tác động tiêu cực đến khả năng sinh lời dài hạn của công ty. Thêm vào đó, các hành động pháp lý phát sinh do vụ vi phạm cũng sẽ gây ra nghĩa vụ tài chính liên tục đáng kể, đặc biệt là nếu các vụ kiện thành công hoặc tiếp tục kéo dài. Do vậy, chúng tôi đưa ra giả thuyết:

\textbf{Giả thuyết 2 (H2)}. \emph{Việc công bố vụ vi phạm dữ liệu có tác động tiêu cực đến giá trị thị trường dài hạn của công ty bị vi phạm.}

Mức độ nghiêm trọng tiềm tàng của vụ vi phạm có khả năng tương quan với số lượng hồ sơ bị lộ. Tài liệu tham khảo [7] đã phát hiện một mối quan hệ thuận chiều giữa số lượng thẻ tín dụng bị lộ và độ lớn của các tỷ suất sinh lợi bất thường âm trong một mẫu rất nhỏ gồm bốn trường hợp. Các phát hiện từ một cuộc khảo sát do [41] thực hiện cũng cho thấy có một mối tương quan thuận giữa số lượng hồ sơ bị lộ và chi phí của vụ vi phạm dữ liệu. Vì vậy, chúng tôi đưa ra giả thuyết:

\textbf{Giả thuyết 3.1 (H3.1)}. \emph{Quy mô của vụ vi phạm dữ liệu có tương quan thuận với tỷ suất sinh lợi âm cao hơn đối với giá trị thị trường ngắn hạn của công ty bị vi phạm.}

Khi một sự kiện vi phạm dữ liệu xảy ra, chúng ta không thể ngay lập tức biết được quy mô chính xác về tác động của sự cố nếu chỉ dựa vào kích thước của vụ rò rỉ. Ví dụ, trong sự cố vi phạm dữ liệu của SONY, sau thông báo đầu tiên về vụ vi phạm, quy mô của các vụ vi phạm tiếp tục mở rộng theo thời gian ngay cả sau khi được phát hiện, do sự tiến hóa trong loại hình xâm nhập của kẻ tấn công và hành vi đánh cắp cơ hội của những kẻ xâm nhập bằng cách ẩn nấp trong các chương trình hoặc các bộ phận khác thuộc hệ thống thông tin của tổ chức. Vì thế, nếu chỉ điều tra mối quan hệ giữa quy mô vụ vi phạm dữ liệu mà một công ty phải gánh chịu và các hiệu ứng ngắn hạn của nó thì không thể là một thước đo hoàn chỉnh về tác động của sự cố. Hơn nữa, các vụ vi phạm dữ liệu lớn hơn sẽ dẫn đến việc mất nhiều khách hàng hơn và các trách nhiệm pháp lý theo sau tương đối nặng nề hơn, và do đó tác động của sự cố cũng sẽ được duy trì trong dài hạn. Vì vậy, không những quy mô của các sự cố không thể được đo lường đầy đủ trong ngắn hạn, mà các thiệt hại pháp lý bồi thường theo sau cũng không thể được đo lường đầy đủ trong giai đoạn ngắn hạn ngay sau sự cố. Ví dụ, trong vụ vi phạm dữ liệu của Target năm 2013, trong đó khối lượng dữ liệu trị giá 70 triệu bảng Anh (pounds) [?] đã bị đánh cắp, một khoản tiền 39 triệu đô la

% [PDF p.7 | in p.7]
đã được chi trả vào năm 2015 cho vụ kiện tập thể [42]. Như chúng ta đã thấy ở đây, mặc dù sự cố vi phạm xảy ra vào năm 2013, khoản tiền bồi thường pháp lý đã không được quyết định cho đến năm 2015. Đã mất ít nhất hai năm để tác động từ sự cố vi phạm của Target được hiện thực hóa đầy đủ. Do đó, chúng tôi đưa ra giả thuyết:

\textbf{Giả thuyết 3.2 (H3.2)}. \emph{Quy mô của vụ vi phạm dữ liệu có tương quan thuận với tỷ suất sinh lợi âm cao hơn đối với giá trị thị trường dài hạn của công ty bị vi phạm.}
